% Este arquivo mostra como inserir uma imagem de fundo no documento através do pacote tikz
\documentclass[a4paper,12pt]{article}
	\usepackage[ansinew]{inputenc}
	\usepackage{ae}
	\usepackage{lipsum}	
	\usepackage[margin=3cm]{geometry}
		
	% O pacote tikz é um conjunto de macros para a construção de figuras. Particularmente, ele permite o posicionamento arbitrário de elementos gráficos no documento, como centralizar uma figura, à imagem do que faz o pacote wallpaper.
	\usepackage{tikz}
	
	\graphicspath{
		{Figuras/EPS/}
		{Figuras/JPG/}
		{Figuras/PDF/}
		{Figuras/PNG/}
		}
		
\begin{document}

	\begin{tikzpicture}[remember picture,overlay]
		\node [opacity=0.50] at (current page.center) {\includegraphics[height=\paperheight,width=\paperwidth]{Tei-Gi}};
	\end{tikzpicture}
	
	\lipsum[1-10]
	
	% Note que o ambiente acima equivale ao comando \ThisCenterWallPaper, haja vista que insere a figura apenas na página que contém o ambiente tikzpicture. Além disso, cabe ressaltar que ao efetuar posicionamento arbitrário de elementos gráficos, como acima, frequentemente são necessárias duas compilações para obter o efeito desejado.
	
\end{document}
